\clearpage
\appendix
\renewcommand{\thesection}{\Alph{section}.}

\section{Hình bổ sung: \texorpdfstring{$\log_{10}\rho$}{log10 rho} theo độ sâu tại bước khởi tạo}
\label{app:phu-luc-a}

Hình~\ref{fig:gradnorm} vẽ lại các cột $\log_{10}\rho$ của Bảng~\ref{tab:grad-init}
dưới dạng đồ thị; dữ liệu trùng với bảng nên hình được đưa về phụ lục.

\begin{figure}[H]
\centering
\begin{tikzpicture}
\begin{axis}[
    width=0.80\linewidth, height=6.0cm,
    xlabel={Độ sâu $L$},
    ylabel={$\log_{10}\rho$ (trung bình $\pm$ độ lệch chuẩn)},
    xtick={5,10,20}, xmin=3, xmax=22,
    ymin=-14, ymax=1,
    axis lines=left,
    legend style={font=\scriptsize, at={(0.03,0.03)}, anchor=south west, draw=none, fill=none},
    tick label style={font=\scriptsize},
    label style={font=\scriptsize},
    grid=major, grid style={gray!20},
]
\addplot[thick, aivnGreen, mark=*, mark size=1.6pt,
         error bars/.cd, y dir=both, y explicit] coordinates {
  (5,-2.78) +- (0,0.358) (10,-6.19) +- (0,0.373) (20,-12.7) +- (0,0.29)
};
\addlegendentry{A: sigmoid + Xavier}
\addplot[thick, brandIndigo, mark=square*, mark size=1.5pt,
         error bars/.cd, y dir=both, y explicit] coordinates {
  (5,-0.74) +- (0,0.135) (10,-0.638) +- (0,0.358) (20,-0.53) +- (0,0.409)
};
\addlegendentry{B: ReLU + He}
\end{axis}
\end{tikzpicture}

\begin{figurecaptionmodern}
\capfig{$\log_{10}\rho$, với $\rho=\|\mathbf{g}^{[1]}\|_F/\|\mathbf{g}^{[L]}\|_F$, theo độ
sâu $L\in\{5,10,20\}$ của MLP rộng $64$ tại bước khởi tạo, trung bình $\pm$ độ lệch
chuẩn trên $5$ hạt giống. Điều kiện A (sigmoid + Xavier) giảm gần tuyến tính theo
$L$ (khoảng $0.2$ mỗi tầng); điều kiện B (ReLU + He) nằm gần $0$ nhưng thanh sai
số rộng dần theo $L$. Trục $y$ là $\log_{10}\rho$ (tương đương thang
log).\label{fig:gradnorm}}
\end{figurecaptionmodern}
\end{figure}
